%=======================================================================
%  CHAPTER 4 - DISCRETE FILTER STRUCTURES  (8 hrs)
%  Every reflection and ladder coefficient quoted here is produced by
%  lattice.py, in exact rational arithmetic, and every quantization
%  answer by quant.py. Both self-test before they print.
%  Numericals live in ch4-num.tex.
%=======================================================================
\section{Discrete Filter Structures}

{\color{sub}\footnotesize 8 hrs \gap$\cdot$\gap \mk{68 questions across the 45 papers}
\gap$\cdot$\gap worth about \mk{10 of 80 marks} \gap$\cdot$\gap 3 syllabus sub-topics
\gap$\cdot$\gap \mk{5th of the 8 chapters} by weight, and tied with Ch 1 as the longest
at 8 hours.}

\lead{Read this first:}
\begin{itemize}
  \item \textbf{42 of the 68 questions are the lattice} \tS{42}, worth 314 of the
        chapter's 449 marks. It is one calculation, the \textbf{step-down recursion},
        run on whatever polynomial the paper prints. Learn it once and two thirds of
        the chapter is finished.
  \item The other two topics are \textbf{realization structures} \tS{16} (direct form I
        and II, cascade, parallel) and \textbf{quantization} \tS{10} (number formats,
        round-off, limit cycles).
  \item \mk{Nothing in this chapter changes what the filter does.} Every structure here
        realizes the \emph{same} $H(z)$. They differ in how much memory they need, how
        many multiplies, and how badly finite word length hurts them. That sentence is
        the answer to "why do we need different structures".
  \item \mk{The lattice is stable exactly when every $|k_m|<1$.} That single line
        answers the "check stability" half of a dozen questions, and it needs no
        root finding at all.
  \item A trap worth knowing before you start: a \textbf{symmetric} polynomial gives
        $|k_N|=1$, and the recursion divides by $1-k_N^2$. \mk{A linear-phase FIR filter
        has no lattice realization.} Five papers ask for one anyway; see \S4.4.
\end{itemize}

\hr

\T{4.1 \quad Why there is more than one structure}

\Q{Why do we need different realization structures for the same system?}{\m{2}\m{4}}
{\color{sub}\footnotesize the opening paragraph for any structure question}

\sbs{%
\lead{What is the same}
\begin{itemize}
  \item Every structure in this chapter realizes the \textbf{same $H(z)$}, so it has the
        same impulse response, the same poles and zeros, and the same frequency
        response. With infinite precision they are indistinguishable.
  \item So a question that asks you to "realize" a system is asking for a \textbf{block
        diagram or signal flow graph}, not for new mathematics.
\end{itemize}}{%
\lead{What differs, and why it matters}
\begin{itemize}
  \item \textbf{Memory}: how many delay elements. Direct Form II uses the minimum
        possible, $\max(M,N)$.
  \item \textbf{Arithmetic}: how many multiplies and adds per output sample.
  \item \textbf{Finite word length}: how far the poles move when the coefficients are
        rounded to $b$ bits, and how much round-off noise accumulates. \mk{This is the
        real reason} cascade and parallel forms are preferred in practice: a direct form
        of high order is extremely sensitive, because every pole depends on every
        coefficient.
  \item \textbf{Modularity}: cascade and lattice are built from identical sections.
\end{itemize}}

\lead{The three building blocks, and the two ways to draw them}
\begin{itemize}
  \item Every structure is made of \textbf{unit delays} $z^{-1}$, \textbf{multipliers}
        (a constant on an arrow) and \textbf{adders}.
  \item A \textbf{block diagram} draws adders as circles with $+$; a \textbf{signal flow
        graph} draws nodes and directed branches, with the gain written on the branch.
        They carry the same information, and papers ask for either.
  \item \textbf{Transposition}: reverse every arrow, swap adders with branch nodes, and
        exchange input with output. The transposed structure has the \emph{same} $H(z)$.
        \mk{No paper in the 45 asks for a transposed form}, and no source in the subject
        covers it, but it is worth one line because it explains how two diagrams that
        look nothing alike can be the same filter.
\end{itemize}

\hr

\T{4.2 \quad Direct Form I and Direct Form II}

\Q{Obtain the Direct Form I and Direct Form II realization of the given system}{\m{2+2}\m{4}\m{5}}
{\color{sub}\footnotesize \tS{16} \yr{\textbf{79 Bh} \m{2+2}, 79 Ba \m{4},
\textbf{82 Bh} \m{2+2}, \textbf{74 Ch} \m{5}, \textbf{72 Ch} \m{5}, 72 Ka \m{4},
\textbf{69 Ch} \m{4}, \textbf{68 Bh} \m{9}, \textbf{70 Ch} \m{4}}
\gap$\cdot$\gap worked in \S1 of the numerical companion}

\sbsr{0.5}{%
\lead{Direct Form I: the difference equation, drawn}
Start from
\[ y[n]=\sum_{k=0}^{M}b_k x[n-k]-\sum_{k=1}^{N}a_k y[n-k] \]
\begin{itemize}
  \item Draw the \textbf{feedforward} part first: a chain of $M$ delays down the input,
        each tapped by $b_k$ into a summing node.
  \item Then the \textbf{feedback} part: a chain of $N$ delays down the output, each
        tapped by $-a_k$ into the same node.
  \item \mk{It is the equation copied onto paper}, which is why it is the easy 2 marks.
  \item Cost: $M+N$ delays, $M+N+1$ multiplies.
  \item Read it as $H(z)=B(z)\cdot\dfrac{1}{A(z)}$: all-zero section first, all-pole
        section second.
\end{itemize}}{%
\lead{Direct Form II: swap the halves and share the delays}
Because the two sections are LTI, they commute:
\[ H(z)=\frac{1}{A(z)}\cdot B(z) \]
\begin{itemize}
  \item Put the \textbf{all-pole section first}. Its output is an intermediate signal
        $w[n]$:
  \item[] \[ w[n]=x[n]-\sum_{k=1}^{N}a_k w[n-k],\qquad
             y[n]=\sum_{k=0}^{M}b_k w[n-k] \]
  \item Now \emph{both} chains delay the same signal $w[n]$, so \mk{one chain of delays
        serves both}.
  \item Cost: $\max(M,N)$ delays, which is the fewest any structure can use. A structure
        with the minimum number of delays is called \textbf{canonic}.
  \item \mk{This is the only reason Direct Form II exists}: same arithmetic, half the
        memory.
\end{itemize}}

\lead{Drawing it, in four steps}
\begin{enumerate}
  \item Write the equation with $y[n]$ alone on the left and \textbf{watch the signs}:
        every $a_k$ crosses over and becomes $-a_k$.
  \item Count $N$ (highest $y$ delay) and $M$ (highest $x$ delay). The chain has
        $\max(M,N)$ delays.
  \item Down the left of the chain put $-a_1,-a_2,\dots$ feeding back into the top
        adder. Down the right put $b_0,b_1,\dots$ feeding forward into the output adder.
  \item Label $w[n]$ at the top of the chain. A missing coefficient is a \textbf{missing
        branch}, not a zero multiplier: if the paper's equation skips $x[n-1]$, leave
        that tap out.
\end{enumerate}

{\color{sub}\footnotesize If $a_0\ne1$, \mk{divide the whole equation by $a_0$ first}.
74 Ch prints $3y[n]+y[n-1]+2y[n-4]=2x[n]+x[n-3]$, and every coefficient must be divided
by 3 before drawing. Forgetting this is the most common lost mark in \S4.2.}

\hr

\T{4.3 \quad Cascade and parallel forms}

\Q{Realize the given $H(z)$ in cascade form of second-order sections, and in parallel
form}{\m{4}\m{5}\m{6}\m{7}}
{\color{sub}\footnotesize \tS{16} \yr{81 Ba \m{5}, \texttt{74 Ma} \m{7}, 80 Ba \m{4},
\textbf{76 Ch} \m{4}, \textbf{\texttt{77 Ch}} \m{7}, \textbf{69 Bh} \m{6},
\textbf{67 Mng}, \textbf{\texttt{78 Ch}} \m{6}}}

\sbsr{0.5}{%
\lead{Cascade: multiply the sections}
\[ H(z)=b_0\prod_{k=1}^{K}\frac{1+b_{1k}z^{-1}+b_{2k}z^{-2}}
                                {1+a_{1k}z^{-1}+a_{2k}z^{-2}} \]
\begin{itemize}
  \item Group the poles and zeros into \textbf{second-order sections}, each realized in
        Direct Form II, and chain them output to input.
  \item \mk{A complex pole must be paired with its conjugate} in the same section. That
        is the whole point: the section then has \textbf{real} coefficients,
        $a_{1k}=-2\mathrm{Re}(p)$ and $a_{2k}=|p|^2$.
\end{itemize}}{%
\lead{Parallel: add the sections}
\[ H(z)=C+\sum_{k=1}^{K}\frac{\gamma_{0k}+\gamma_{1k}z^{-1}}
                              {1+a_{1k}z^{-1}+a_{2k}z^{-2}} \]
\begin{itemize}
  \item Get here by \textbf{partial fractions}, exactly as in Chapter 2, then recombine
        each conjugate pair into one real second-order section.
  \item $C$ is the constant from the division when $H(z)$ is improper; it is a direct
        path from input to output.
\end{itemize}}

\sbsr{0.5}{%
\begin{itemize}
  \item Real poles pair with real poles, or sit alone as a first-order section.
  \item An odd number of real poles leaves one first-order section; that is normal, not
        an error.
\end{itemize}}{%
\begin{itemize}
  \item \mk{Cascade needs factors, parallel needs partial fractions.} That is the only
        difference in the work.
  \item Parallel has the best round-off behaviour of the three: an error in one section
        cannot propagate into the others, because they do not feed each other.
\end{itemize}}

\lead{Turning a conjugate pair into real coefficients}
\begin{itemize}
  \item A pair of poles at $re^{\pm j\theta}$ multiplies out to
  \item[] \[ (1-re^{j\theta}z^{-1})(1-re^{-j\theta}z^{-1})
             =1-2r\cos\theta\,z^{-1}+r^{2}z^{-2} \]
  \item So a section written $1-2r\cos\theta\,z^{-1}+r^{2}z^{-2}$ needs only $r$ and
        $\theta$, both of which the papers usually hand you directly.
  \item For a pair given in rectangular form $\sigma\pm j\omega$, it is
        $1-2\sigma z^{-1}+(\sigma^{2}+\omega^{2})z^{-2}$.
  \item \mk{Check every section has real coefficients before drawing.} A surviving $j$
        means a pole was left unpaired.
\end{itemize}

\hr

\T{4.4 \quad The FIR lattice}

\Q{Compute the lattice coefficients and draw the structure for the given $H(z)$; check
stability}{\m{5}\m{6}\m{8}\m{10}}
{\color{sub}\footnotesize \tS{42} \yr{\textbf{80 Bh} \m{10},
\textbf{\texttt{81 Ch}}, \textbf{\texttt{72 Ash}} \m{4+2+1}, \textbf{\texttt{79 Ch}},
\textbf{\texttt{69 Bh}} \m{8}, \textbf{\texttt{71 Bh}} \m{6+2}, \texttt{70 Ma} \m{6},
\textbf{71 Ch} \m{4+2+1}, 81 Ba \m{5}, \textbf{78 Bh} \m{3+7}, \textbf{79 Bh} \m{6},
71 Shr \m{5}, \textbf{\texttt{80 Ch}}, \texttt{74 Ma} \m{7}, and 25 more}
\gap$\cdot$\gap \mk{the single most valuable calculation in the subject}}

\sbsr{0.52}{%
\lead{What the structure is}
An $M$-stage lattice chains identical two-input two-output sections. Stage $m$ takes a
forward signal $f_{m-1}$ and a backward signal $g_{m-1}$ and produces
\[ f_m[n]=f_{m-1}[n]+k_m\,g_{m-1}[n-1] \]
\[ g_m[n]=k_m\,f_{m-1}[n]+g_{m-1}[n-1] \]
with $f_0[n]=g_0[n]=x[n]$ and $y[n]=f_M[n]$.
\begin{itemize}
  \item The $k_m$ are the \textbf{reflection coefficients}, one per stage.
  \item The forward path realizes $A_M(z)$; \mk{the backward path realizes its reverse}
        $z^{-M}A_M(z^{-1})$, which is what makes the all-pass of \S4.5 so tidy.
  \item Only \textbf{one} multiplier value per stage, used twice. That symmetry is why
        the lattice is far less sensitive to coefficient quantization than a direct
        form, and why it is used for speech.
\end{itemize}}{%
\lead{The step-down recursion \mk{(direct form $\rightarrow$ lattice)}}
Given $A_M(z)=1+a_M(1)z^{-1}+\dots+a_M(M)z^{-M}$:
\[ k_m=a_m(m) \]
\[ a_{m-1}(i)=\frac{a_m(i)-k_m\,a_m(m-i)}{1-k_m^{2}},\quad i=1,\dots,m-1 \]
\begin{enumerate}
  \item \textbf{Read off $k_M$}: it is the last coefficient, free.
  \item Fold $A_M$ against its own reverse, divide by $1-k_M^{2}$, and you have
        $A_{M-1}$, one degree lower.
  \item Repeat down to $A_1$, whose single coefficient is $k_1$.
\end{enumerate}
\lead{The step-up recursion \mk{(lattice $\rightarrow$ direct form)}}
\[ a_m(i)=a_{m-1}(i)+k_m\,a_{m-1}(m-i),\qquad a_m(m)=k_m \]
{\color{sub}\footnotesize Same relation, rearranged. Papers ask in both directions:
79 Bh and 71 Shr give $k=\{\tfrac14,\tfrac12,\tfrac13\}$ and want the FIR coefficients.}}

\lead{Stability, in one line}
\begin{itemize}
  \item \mk{The system is stable if and only if $|k_m|<1$ for every $m$.}
  \item No root finding, no Jury table. The step-down \emph{is} the Schur-Cohn test: it
        is testing whether every root of $A(z)$ lies inside the unit circle.
  \item A question that says "check stability" after asking for the lattice is asking
        you to \textbf{look at the numbers you already computed}. It is a free mark.
\end{itemize}

\creamq{The symmetric trap: why five papers ask for an impossible lattice}{If $A(z)$ is
\textbf{symmetric}, meaning its coefficient list reads the same backwards, then
$a_M(M)=a_M(0)=1$, so $k_M=1$ and the step-down divides by $1-k_M^{2}=0$. \mk{A
symmetric polynomial has no lattice realization.} And symmetry is exactly the
linear-phase condition $h[n]=h[M-n]$ from Chapter 3, so \textbf{a linear-phase FIR
filter can never be put in lattice form}. The papers that hit this are
\textbf{72 Ch} ($1+2z^{-1}+z^{-2}$), \textbf{\texttt{73 Bh}} ($1+2z^{-1}+2z^{-2}+z^{-3}$),
\textbf{75 Ch} and \textbf{73 Ch} (both with $1+\tfrac23z^{-1}+\tfrac58z^{-2}
+\tfrac23z^{-3}+z^{-4}$), and 74 Ash, whose last coefficient is $-1$. Say so, show the
division by zero, and note that the zeros sit \emph{on} the unit circle. It is a better
answer than a page of arithmetic that cannot finish.}

\hr

\T{4.5 \quad Lattice-ladder, for a pole-zero system}

\sbsr{0.5}{%
\lead{Why a ladder is needed}
\begin{itemize}
  \item The plain lattice realizes an \textbf{all-zero} $A_M(z)$, or upside down an
        \textbf{all-pole} $1/A_N(z)$.
  \item A general $H(z)=B(z)/A(z)$ needs one more piece: tap every backward signal
        $g_m[n]$ through a weight $C_m$ and sum them.
  \item[] \[ y[n]=\sum_{m=0}^{M}C_m\,g_m[n] \]
  \item The lattice supplies the poles, \mk{the ladder taps supply the zeros}.
\end{itemize}}{%
\lead{Getting the ladder coefficients}
Run the step-down on the \textbf{denominator} $A(z)$ and keep every intermediate
polynomial $A_i(z)$. Then, working \textbf{downwards} from $m=M$:
\[ C_m=b_m-\sum_{i=m+1}^{M}C_i\,a_i(i-m) \]
\begin{enumerate}
  \item $C_M=b_M$, because the sum is empty.
  \item Each lower $C$ subtracts what the higher taps already contribute.
  \item \mk{The $k$'s come only from $A(z)$}, so two papers with the same denominator
        have the same lattice and differ only in the ladder. 76 Ch and 76 Ash are such a
        pair, and so are 74 Bh and 75 Bh.
\end{enumerate}}

\creamq{The all-pass result, and why it is worth remembering}{\textbf{\texttt{70 Bh}}
gives $H(z)=\dfrac{1+0.4z^{-1}-1.2z^{-2}+2z^{-3}}{2-1.2z^{-1}+0.4z^{-2}+z^{-3}}$ and asks
how a lattice implements it. Look at the coefficients: \mk{the numerator is the
denominator reversed}. That is the definition of an all-pass,
$H(z)=z^{-N}A(z^{-1})/A(z)$, and the backward path of a lattice already computes exactly
$z^{-N}A(z^{-1})$. So the ladder collapses: every $C_m$ is zero except $C_N=1$. \mk{An
all-pass filter is just a lattice with the output taken from the last backward node},
no ladder taps at all. The reflection coefficients are $k=\{-0.56,\ \tfrac23,\
\tfrac12\}$, all less than 1 in magnitude, so it is stable.}

\hr

\T{4.6 \quad Number representation and quantization}

\Q{Write about the sign magnitude and 2's complement representation of a binary
fractional number. Write about truncation error and rounding error}{\m{6}\m{3+4}}
{\color{sub}\footnotesize \tS{10} \yr{\textbf{69 Bh} \m{6}, \textbf{\texttt{69 Bh}} \m{3+4},
75 Ash \m{3}, \textbf{67 Mng}, 71 Shr \m{1+1}}}

\sbs{%
\lead{The three fixed-point codes}
A fraction in $[-1,1)$ with $b$ bits after the binary point. Positive numbers look the
same in all three; only negatives differ.
\par\smallskip
\begin{tabular}{@{}L{2.5cm}L{5.0cm}@{}}
\toprule
\hrow \thead{Code} & \thead{How a negative is formed} \\
Sign magnitude & sign bit $1$, magnitude bits unchanged \\
1's complement & invert every magnitude bit \\
2's complement & invert every bit, add $1$ in the last place \\
\bottomrule
\end{tabular}
\par\smallskip
\begin{itemize}
  \item Worked, for $\pm5/8$ in 3 bits, $5/8=0.101_2$:
  \item[] $+5/8=\mathtt{0.101}$ in all three.
  \item[] $-5/8=\mathtt{1.101}$ (SM), $\mathtt{1.010}$ (1's), $\mathtt{1.011}$ (2's).
  \item \mk{2's complement is used in practice} because addition needs no special case
        and overflows in a partial sum cancel out.
\end{itemize}}{%
\lead{Rounding against truncation}
With step $q=2^{-b}$ and error $e=Q(x)-x$:
\par\smallskip
\begin{tabular}{@{}L{3.2cm}L{4.3cm}@{}}
\toprule
\hrow \thead{Scheme} & \thead{Error range} \\
Rounding, any code & $-q/2\le e\le q/2$ \\
Truncation, sign mag. & $-(q-q_u)\le e\le(q-q_u)$ \\
Truncation, 2's comp. & $-(q-q_u)\le e\le0$ \\
\bottomrule
\end{tabular}
\par\smallskip
\begin{itemize}
  \item \mk{Rounding is better}, and this table is why: its error is centred on zero, so
        it has \textbf{no bias}, and its range is half as wide.
  \item Truncating a \textbf{sign magnitude} number always shrinks the magnitude, so the
        error is symmetric about zero but takes the sign of $x$.
  \item Truncating a \textbf{2's complement} number always moves the value \emph{down},
        never up, so the error is \textbf{one-sided and negative}: a systematic bias, a
        DC offset that accumulates through a filter.
  \item $q_u=2^{-b_u}$ is the step before quantization, with $b_u>b$.
\end{itemize}}

\lead{The three places quantization bites}
\begin{itemize}
  \item \textbf{Input (A/D) quantization}: modelled as additive white noise of variance
        $q^{2}/12$ for rounding.
  \item \textbf{Coefficient quantization}: the $a_k$ and $b_k$ are stored in $b$ bits, so
        \mk{the poles move}. A high-order direct form is the worst case, because every
        pole depends on every coefficient; cascade and parallel forms confine the damage
        to one section, which is the practical argument for \S4.3.
  \item \textbf{Product round-off}: each multiply doubles the word length and must be
        shortened again. In a recursive filter this feeds back, which is \S4.7.
\end{itemize}

\hr

\T{4.7 \quad Limit cycles and the dead band}

\Q{What is the limit cycle effect in a recursive system? Describe it with an example.
Define dead band}{\m{3}\m{5}\m{1+3}}
{\color{sub}\footnotesize \tS{10} \yr{\textbf{71 Ch} \m{3}, 70 Asa \m{1+3},
\texttt{72 Ma} \m{5}, \textbf{\texttt{81 Ch}}, \textbf{\texttt{80 Ch}},
\textbf{\texttt{70 Bh}} \m{3+4}, \textbf{66 Ma} \m{4}, 71 Shr \m{1+1+2+1}}}

\sbsr{0.5}{%
\lead{What a limit cycle is}
\begin{itemize}
  \item In a \textbf{recursive} filter the rounded product is fed back and rounded again.
        Once the signal is small enough, rounding can return the \emph{same magnitude} it
        started from.
  \item The filter then never decays: after the input stops it sits in a
        \textbf{self-sustaining oscillation} of fixed amplitude. That is a
        \textbf{zero-input limit cycle}.
  \item It happens only with quantization. \mk{The ideal infinite-precision filter is
        perfectly stable}; the limit cycle is created by the arithmetic, not by the poles.
\end{itemize}}{%
\lead{The dead band}
For the first-order $y[n]=Q\{a\,y[n-1]\}+x[n]$ with rounding, the filter cannot escape
once
\[ |y[n]|\le\frac{q}{2\,(1-|a|)} \]
\begin{itemize}
  \item That interval is the \textbf{dead band}: inside it, quantization exactly undoes
        the decay that $a$ should have produced.
  \item The oscillation amplitude is bounded by the dead band, so a longer word (smaller
        $q$) makes it smaller but never removes it.
\end{itemize}}

\sbsr{0.5}{%
\begin{itemize}
  \item FIR filters cannot have them: no feedback, nothing to sustain.
\end{itemize}}{%
\begin{itemize}
  \item If $a>0$ the trapped value holds \textbf{constant}, period 1. If $a<0$ it
        \textbf{alternates in sign}, period 2.
\end{itemize}}

\lead{The worked example, 66 Ma \mk{(and the shape of every limit-cycle answer)}}
$y[n]=Q\{-0.5\,y[n-1]\}+x[n]$, rounded to 3 bits so $q=\tfrac18$, with $y[-1]=0$ and
$x[n]=0.875\,\delta[n]$:
\par\smallskip
\begin{tabular}{@{}L{1.0cm}L{2.0cm}L{2.6cm}L{2.6cm}L{4.0cm}@{}}
\toprule
\hrow \thead{$n$} & \thead{$-0.5\,y[n-1]$} & \thead{rounded to $\tfrac18$}
& \thead{$+\,x[n]$} & \thead{} \\
0 & $0$ & $0$ & $0.875$ & the impulse arrives \\
1 & $-0.4375$ & $-0.5$ & $-0.5$ & \\
2 & $0.25$ & $0.25$ & $0.25$ & \\
3 & $-0.125$ & $-0.125$ & $-0.125$ & now inside the dead band \\
4 & $0.0625$ & $\mathbf{0.125}$ & $\mathbf{0.125}$ & \mk{rounding put it back up} \\
5 & $-0.0625$ & $\mathbf{-0.125}$ & $\mathbf{-0.125}$ & \\
\bottomrule
\end{tabular}
\par\smallskip
\begin{itemize}
  \item From $n=3$ on it alternates $\pm\tfrac18$ for ever: a limit cycle of amplitude
        $\tfrac18$ and \textbf{period 2}, the period coming from $a<0$.
  \item Dead band: $\dfrac{q}{2(1-|a|)}=\dfrac{1/8}{2(1-0.5)}=\dfrac18$, which is exactly
        the amplitude observed. \mk{Quote the formula and the table agrees with it.}
  \item Step 4 is the whole mechanism in one line: $-0.5\times(-0.125)=0.0625$, which is
        \emph{half} a step, and rounding carries it up to a full $0.125$ instead of down
        to $0$.
\end{itemize}

\hr

\T{4.8 \quad Last-minute recall}

\sbs{%
\lead{Say these without thinking}
\begin{itemize}
  \item Direct Form II is \textbf{canonic}: $\max(M,N)$ delays, the fewest possible.
  \item Divide through by $a_0$ before drawing anything.
  \item Cascade needs \textbf{factors}; parallel needs \textbf{partial fractions}.
  \item A complex pole always sits with its conjugate, giving
        $1-2r\cos\theta\,z^{-1}+r^{2}z^{-2}$.
  \item Step-down: $k_m=a_m(m)$, then
        $a_{m-1}(i)=\dfrac{a_m(i)-k_m a_m(m-i)}{1-k_m^{2}}$.
  \item Step-up: $a_m(i)=a_{m-1}(i)+k_m a_{m-1}(m-i)$.
  \item \textbf{Stable} $\Leftrightarrow$ every $|k_m|<1$.
  \item Ladder: $C_m=b_m-\sum_{i>m}C_i\,a_i(i-m)$, downwards from $C_M=b_M$.
  \item Dead band $=\dfrac{q}{2(1-|a|)}$; $a<0$ gives period 2.
\end{itemize}}{%
\lead{Mistakes that cost marks}
\begin{itemize}
  \item \mk{Forgetting to divide by $a_0$} when the paper prints $3y[n]+\dots$.
  \item Sign errors moving $a_k$ across the equals sign: the feedback taps are
        $\mathbf{-a_k}$.
  \item Running the step-down on the \textbf{numerator} of a lattice-ladder question.
        The $k$'s come from $A(z)$, always.
  \item Grinding through a step-down that has already produced $|k|=1$. Stop and say the
        polynomial is symmetric.
  \item Reporting $k$ to two decimals and then calling $|k|=0.99$ stable without saying
        so. Keep fractions where the paper gives fractions.
  \item Claiming an FIR filter can have a limit cycle. It cannot: no feedback.
  \item Saying truncation error is symmetric for 2's complement. \mk{It is one-sided},
        and that is the point of the question.
\end{itemize}}
